\chapter{Scout Agent}

Lo Scout \`e ottimizzato per la copertura massima della griglia nel minor numero di passi.
Non recupera oggetti e non riceve compiti: la sua unica responsabilit\`a \`e esplorare
e condividere informazioni.

\section{Ciclo Decisionale}

La priorit\`a dello Scout \`e deterministica e segue questa sequenza:

\begin{enumerate}
    \item \textbf{Ricarica energetica}: Se l'energia scende sotto la soglia critica (calcolata
          come la distanza stimata al magazzino pi\`u vicino moltiplicata per un fattore di
          sicurezza), lo Scout interrompe l'esplorazione e si dirige al magazzino. Durante
          la ricarica, continua a comunicare la mappa locale e gli oggetti scoperti.
    \item \textbf{Consegna scoperte al Coordinator}: Se lo Scout ha scoperto nuovi oggetti
          e un Coordinator si trova nel raggio di comunicazione, li trasmette immediatamente
          tramite \texttt{ObjectSpottedMessage}. Se nessun Coordinator \`e a portata, lo Scout
          si dirige verso l'ultima posizione nota del Coordinator (\texttt{last\_seen\_coordinator\_pos})
          per consegnare le informazioni al pi\`u presto. Se \`e trascorso troppo tempo senza
          trovarlo, viene effettuata una scansione estesa fino a $4\times$ il raggio di
          comunicazione per localizzarlo.
    \item \textbf{Selezione frontiera}: Se non \`e in ricarica, \texttt{FrontierExplorer}
          calcola le frontiere disponibili dalla mappa locale e seleziona quella con utilit\`a
          massima. La funzione $U(f)$ favorisce cluster grandi e vicini, penalizzando cluster
          gi\`a puntati da altri Scout (comunicazione implicita via \texttt{target\_position}
          osservata nei messaggi di stato).
    \item \textbf{Movimento}: A* calcola il percorso ottimale verso la frontiera selezionata,
          con le posizioni degli altri agenti come penalizzazione soft.
    \item \textbf{Comunicazione}: Ad ogni passo, lo Scout condivide con tutti gli agenti
          entro il raggio \texttt{comm\_radius} la propria mappa locale aggiornata tramite
          \texttt{MapDataMessage} e la lista degli oggetti scoperti tramite
          \texttt{ObjectSpottedMessage} indirizzati ai Coordinator vicini.
\end{enumerate}

\section{Selezione delle Frontiere}

\subsection{Funzione di Utilit\`a con Distanza Sub-lineare}

La funzione di utilit\`a per la selezione delle frontiere utilizza una distanza con
esponente sub-lineare:

$$U(f) = \frac{\text{cluster\_size}}{d(f)^{0.4} + 1} - 0{,}5 \cdot n_{\text{agenti}}(f)$$

L'uso di $d^{0.4}$ anzich\'e $d$ lineare impedisce che cluster piccoli ma vicini
prevalgano sistematicamente su cluster grandi e distanti. Senza questa correzione,
un cluster di 3~celle a distanza~2 vincerebbe contro un cluster di 50~celle a
distanza~25, portando a una copertura frammentata.

\subsection{Blocco del Target (Target Lock)}

Quando una frontiera viene selezionata, lo Scout imposta un \textit{target lock}
per un numero configurabile di passi (\texttt{\_TARGET\_LOCK\_DURATION}). Durante
questo periodo, la rivalutazione completa delle frontiere \`e sospesa anche se ne
compaiono di nuove. Il lock viene rilasciato anticipatamente se il target \`e
raggiunto, se il rilevatore di stallo si attiva, o se la durata scade. Questo
meccanismo previene cambi di direzione erratici quando lo Scout entra in zone con
molti cluster di frontiera adiacenti.

\subsection{Blacklist dei Target Recenti}

Dopo aver raggiunto una frontiera, la sua posizione viene inserita in una blacklist
per \texttt{\_RECENT\_TARGET\_TTL} passi (default 200). Ci\`o previene l'oscillazione
immediata tra cluster appena esplorati: senza la blacklist, un cluster marginalmente
non coperto verrebbe continuamente riselezionato prima di essere completamente
scansionato.

\subsection{Preferenza per Frontiere Lontane}

Quando abilitata (\texttt{\_FAR\_FRONTIER\_ENABLED}), la selezione preferisce frontiere
che richiedono uno spostamento superiore a $2\times$ il raggio di visione (distanza minima
configurabile, default 8~celle). Questo spinge lo Scout verso territorio genuinamente
inesplorato, evitando di micro-esplorare artefatti di confine tra cluster adiacenti gi\`a
parzialmente coperti. Se non esistono frontiere lontane, si ricade sulle frontiere
disponibili. In modalit\`a \texttt{map\_known}, questa euristica viene disattivata:
le frontiere vicine al confine tra zone scansionate e non scansionate sono le pi\`u
rilevanti.

\section{Anti-Clustering}

L'anti-clustering \`e implementato in due livelli. Il primo \`e un filtro hard: le frontiere
che si trovano a meno di \texttt{\_ANTI\_CLUSTER\_DIST} celle (default 8) da un altro Scout
(visibile nel raggio di comunicazione) vengono escluse completamente dall'insieme dei candidati.
Il secondo \`e una penalizzazione soft nella funzione di utilit\`a:

$$U(f) = \frac{\text{cluster\_size}}{d(f)^{0.4} + 1} - 0{,}5 \cdot n_{\text{agenti}}(f)$$

dove $n_{\text{agenti}}(f)$ \`e il numero di agenti entro 10 celle dalla frontiera $f$.
Se il filtro hard elimina tutte le frontiere (caso raro, ad esempio con molti Scout su
una griglia piccola), si ricade sull'insieme completo applicando solo la penalizzazione soft.
Il risultato pratico \`e una divisione organica del territorio senza alcuna coordinazione esplicita.

\section{Pattugliamento Ciclico (Stale Coverage)}

Ogni Scout mantiene una matrice privata \texttt{\_coverage\_step} che registra il
passo del modello in cui ciascuna cella \`e stata l'ultima volta nel raggio di visione.
Questa matrice non viene mai trasmessa.

Quando non esistono frontiere inesplorate, lo Scout entra in modalit\`a di
\textit{stale coverage patrol}: le celle non visitate per un numero configurabile di
passi (\texttt{\_RESCAN\_AGE}) diventano candidati per la ri-esplorazione.
La selezione combina et\`a e prossimit\`a:

$$\text{score}(c) = \frac{\text{age}(c)}{d(c) + 1}$$

I primi 50 candidati per punteggio vengono valutati; il primo raggiungibile (cella
walkable) viene selezionato come target. In modalit\`a \texttt{map\_known}, le celle
con \texttt{vision\_explored == 0} hanno priorit\`a assoluta rispetto al
ri-pattugliamento di celle gi\`a scansionate.

\section{Relay Passivo e Ricerca del Coordinator}

Quando lo Scout ha oggetti scoperti da consegnare, valuta se la catena di relay
passivo (via \texttt{MapDataMessage} trasmesso ad agenti vicini) \`e sufficiente.
Se il contatto con un qualsiasi agente \`e avvenuto di recente
(\texttt{\_last\_agent\_contact\_step}), lo Scout presume che le informazioni stiano
gi\`a propagandosi e continua a esplorare anzich\'e cercare attivamente il Coordinator.

Se il contatto \`e stale, lo Scout si dirige verso l'ultima posizione nota del
Coordinator. All'arrivo, se il Coordinator non \`e presente, esegue una
\textit{wide scan} fino a $4\times$ il raggio di comunicazione per localizzarlo
senza navigazione cieca.

\section{Finestra di Rilevamento Stallo Raddoppiata}

A differenza degli altri agenti (che usano uno storico di 15 posizioni), lo Scout
utilizza \texttt{max\_history\_length = 30}. Poich\'e gli Scout si muovono a
$2\times$ la velocit\`a degli altri agenti, 15~posizioni coprono solo $\sim$7
passi reali. Il raddoppio garantisce che la finestra di rilevamento loop copra
$\sim$15 passi effettivi, evitando falsi positivi.
